% =============================================================================
% Partie A — Architecture globale
% =============================================================================
\section{Partie A — Architecture globale}

\subsection{Diagramme de contexte}

Le système a trois catégories d'utilisateurs, qui n'accèdent pas au service de la
même façon.

Le \textbf{rédacteur} s'identifie par un nom et une adresse électronique, sans
mot de passe, et reçoit un jeton de session. Il crée des documents, les partage,
les versionne et fabrique des liens de téléchargement.

Le \textbf{visiteur} est le cas intéressant : il ne s'identifie \emph{jamais}.
Le jeton du lien qu'il détient \emph{est} sa capacité d'accès. Cette décision a
une conséquence architecturale : le parcours visiteur est exposé à la racine
(\code{/d/\{token\}/\dots}), en dehors du préfixe \code{/api} qui regroupe les
routes authentifiées. La confondre avec les autres routes produit un échec
silencieux : la requête aboutit sur le fichier statique du front et renvoie un
code 200 avec du HTML.

L'\textbf{administrateur} consulte l'état du système avec une clé dédiée, dans
un en-tête, et non avec le jeton de session : c'est une capacité
d'exploitation, pas une identité applicative.

Deux éléments extérieurs complètent le tableau. \textbf{Amazon SES} envoie un
unique courriel lorsque le dernier exemplaire d'un lien est distribué.
\textbf{Le stockage objet} sert les fichiers PDF : le système ne les transporte
jamais, il présigne une adresse et laisse S3 répondre. C'est ce qui permet à
500 téléchargements de ne pas passer par l'instance applicative.

\illus{01-contexte}{0.82}{Diagramme de contexte : trois acteurs, un système,
deux services extérieurs.}{fig:contexte}

Un point de conception mérite d'être explicité ici. Le destinataire du courriel
de fin de quota n'est pas un acteur du système : il ne dispose d'aucune
interface, ne peut rien demander et ne reçoit rien d'autre. Le représenter comme
un acteur donnerait une idée fausse du périmètre.

\subsection{Architecture de déploiement AWS}

Tout le trafic entre par une \textbf{adresse IP élastique}, attachée à une
instance unique. Cette adresse est indispensable : une adresse publique
ordinaire change à chaque arrêt puis démarrage de l'instance, ce qui rendrait
l'évaluation impossible avec une URL stable.

Sur l'instance cohabitent trois rôles, séparés par des processus distincts et
non par des machines :

\begin{itemize}[leftmargin=1.4em,itemsep=2pt]
  \item \textbf{nginx} sert le front statique, met en micro-cache l'état des
        liens, et relaie vers Node les préfixes \code{/api}, \code{/d} et
        \code{/ws} ;
  \item \textbf{le processus applicatif} tient l'état Yjs des documents ouverts
        en mémoire et sert l'API REST et le WebSocket ;
  \item \textbf{le processus de travail} consomme la file de génération de PDF,
        séparément pour qu'un document lourd ne bloque pas les écritures.
\end{itemize}

\illus{02-deploiement}{0.95}{Architecture de déploiement : une région, une
instance applicative, quatre services gérés.}{fig:deploiement}

Le choix du type d'instance est une décision de dimensionnement, pas de confort.
Le type retenu est \code{m7i-flex.large} : deux processeurs virtuels et huit
gigaoctets de mémoire. Deux raisons, dans cet ordre :

\begin{enumerate}[leftmargin=1.4em,itemsep=2pt]
  \item c'est le seul type à la fois \textbf{éligible à l'offre gratuite} et doté
        de la mémoire nécessaire pour tenir des milliers de connexions
        simultanées. \code{t3.large} a été essayé d'abord : la création de pile
        échoue avec « not eligible for Free Tier » ;
  \item les instances de la famille \code{t} consomment des crédits de processeur
        et se retrouvent bridées quand ces crédits s'épuisent. En mode
        \emph{Standard}, une instance bridée s'effondre précisément pendant le
        pic qu'on cherche à mesurer. Le mode \emph{illimité} évite le
        bridage mais rend la facture imprévisible.
\end{enumerate}

\subsection{Décisions d'architecture et alternatives écartées}

Chaque décision est consignée dans un document d'architecture séparé
(\code{docs/adr/}), avec ses alternatives. Le tableau ci-dessous en donne la
synthèse et, surtout, ce qui a été \emph{écarté} : une décision sans alternative
n'est pas une décision.

\begin{table}[H]
\centering
\small
% Largeur utile = 21 cm − 2×2,4 cm de marge − 4×4 pt d'intercolonne = 15,64 cm.
% Le total ci-dessous (15,5 cm) tient avec une marge ; à 15,9 cm il débordait.
\begin{tabular}{@{}p{3.0cm}p{6.2cm}p{6.3cm}@{}}
\toprule
\textbf{Choix} & \textbf{Retenu, et pourquoi} & \textbf{Écarté, et pourquoi} \\
\midrule
Calcul
& \textbf{Instance EC2 unique} : le WebSocket et l'état en mémoire sont
  naturels, et un seul écrivain par document rend C2 simple à tenir
& \textbf{Lambda + API Gateway} : la limite de 400 exécutions simultanées du
  compte est trop serrée pour 15\,000 clics, et le WebSocket persistant
  obligerait à externaliser tout l'état \\
\addlinespace
Convergence
& \textbf{CRDT Yjs} : les modifications concurrentes fusionnent sans verrou, et
  un client hors ligne se réconcilie à la reconnexion
& \textbf{Verrou exclusif} : un seul rédacteur à la fois, ce que le sujet
  interdit ; \textbf{transformation opérationnelle} : juste, mais coûteux à
  écrire et à prouver dans le temps d'un projet \\
\addlinespace
Écriture
& \textbf{Écrivain unique et lots de 250\,ms} : une transaction par lot au lieu
  d'une par frappe
& \textbf{Écriture par frappe} : une transaction DynamoDB par caractère épuise
  le débit disponible et multiplie la facture \\
\addlinespace
Persistance
& \textbf{DynamoDB} : transactions conditionnelles, ce qui donne C1 exactement,
  et aucune administration de serveur
& \textbf{PostgreSQL} : les transactions seraient plus expressives, mais
  imposent un serveur à administrer, sauvegarder et dimensionner \\
\addlinespace
Fichiers
& \textbf{S3 privé + URL présignée de 60 s} : le fichier ne traverse pas
  l'instance
& \textbf{Flux servi par Node} : 500 PDF simultanés saturent la bande passante
  de l'instance et bloquent l'édition \\
\addlinespace
PDF
& \textbf{Génération asynchrone par file} : le visiteur reçoit « patienter » et
  revient ; un seul travail par version grâce à un verrou
& \textbf{Génération dans la requête} : une requête de plusieurs secondes sous
  15\,000 clics ferait tomber l'instance \\
\addlinespace
Pic visiteur
& \textbf{Micro-cache nginx sur l'état du lien} : une seconde de cache absorbe
  15\,000 lectures sans toucher la base
& \textbf{Interrogation directe de DynamoDB} : 15\,000 lectures par seconde la
  font échouer en limitation de débit \\
\addlinespace
Front
& \textbf{Export statique servi par nginx} : quelques fichiers, aucun serveur
  applicatif pour le rendu
& \textbf{Rendu côté serveur Next.js} : un second processus Node à
  dimensionner, sans bénéfice ici \\
\bottomrule
\end{tabular}
\caption{Décisions d'architecture et alternatives écartées.}
\label{tab:adr}
\end{table}

\subsection{Flux internes}

Deux flux méritent une figure : celui qui tient la contrainte C1, et celui qui
tient C2 avec C4. Ce sont les deux contraintes éliminatoires.

\subsubsection{Réservation d'un exemplaire}

L'ordre des vérifications est imposé par le sujet et a des conséquences
observables. Avant de décider que le lien est épuisé, le serveur cherche
\emph{d'abord} si ce visiteur possède déjà un ticket. Sans cette précaution, un
visiteur déjà servi recevrait « épuisé » après l'épuisement du lien — il perdrait
un exemplaire qu'il avait légitimement obtenu.

Ensuite, la réservation est une transaction conditionnelle : décrémenter
\code{remaining} \emph{si} il est supérieur à zéro et si le lien est actif, et
créer le ticket \emph{si} ce visiteur n'en a pas. La transition vers l'état
épuisé est une écriture séparée, elle aussi conditionnelle : seul le claim qui
amène réellement le compteur à zéro peut la poser. C'est cette séparation qui
rend l'épuisement idempotent lorsque cinquante visiteurs arrivent ensemble.

\illus{03-flux-reservation}{0.92}{Flux de réservation : transaction
conditionnelle, puis génération asynchrone du PDF.}{fig:flux-reservation}

Le numéro de séquence est figé au moment du claim, avant la transaction. Tous
les visiteurs d'une même rafale partagent donc le même fichier, et une
modification ultérieure du document ne change pas leur exemplaire.

\subsubsection{Écriture groupée}

Le flux d'écriture porte deux contraintes à la fois : aucune modification perdue
(C2) et aucun dépassement de crédit (C4). Les deux se règlent au même endroit,
parce qu'elles portent sur la même transaction.

\illus{04-flux-ecriture}{0.92}{Flux d'écriture : un lot par tick de 250\,ms, une
transaction atomique, un seul message diffusé.}{fig:flux-ecriture}

Un point de conception a coûté cher et mérite d'être rapporté. Le contrôle du
crédit ne peut pas être une condition de la transaction : DynamoDB refuse
l'arithmétique entre deux valeurs dans une expression de condition. La forme
\code{used + :chars <= :limit} est rejetée par le service. Le contrôle se fait
donc par une lecture préalable, hors transaction. Le risque de course est borné
par l'écrivain unique ; la transaction reste la garde d'atomicité pour la
séquence, le journal et les marqueurs d'idempotence.

\subsection{Sécurité et bonnes pratiques}

Les mesures ci-dessous répondent directement aux bonnes pratiques du cours.

\begin{table}[H]
\centering
\small
\begin{tabular}{@{}p{4.6cm}p{11cm}@{}}
\toprule
\textbf{Point} & \textbf{Mise en œuvre} \\
\midrule
Aucune clé d'accès dans le code
& le poste utilise une connexion fédérée ; le serveur utilise le rôle de
  l'instance, donc des identifiants temporaires renouvelés \\
\addlinespace
Aucun secret dans les données de démarrage
& les données de démarrage d'une instance sont lisibles via le service de
  métadonnées. Les deux secrets (signature des sessions, clé d'administration)
  sont dans le magasin de paramètres, créés une fois et conservés \\
\addlinespace
Moindre privilège
& le rôle de l'instance n'accède qu'aux tables nommées, au préfixe de stockage,
  à la file, et à l'envoi de courriel filtré sur l'expéditeur vérifié \\
\addlinespace
Métadonnées d'instance
& jeton obligatoire : sans session, les données de démarrage ne sont pas
  lisibles \\
\addlinespace
Stockage objet
& accès public bloqué, chiffrement au repos, versionnage, règle refusant le
  trafic non chiffré, URL présignées de 60 secondes \\
\addlinespace
Pare-feu
& port 80 ouvert ; port 22 \emph{fermé} par défaut, accès par session
  sécurisée. Une valeur permissive est explicitement traitée comme une
  fermeture, et non comme une ouverture \\
\addlinespace
Jeton WebSocket
& le navigateur ne permet pas d'en-tête d'authentification au handshake : un
  ticket à usage unique de 60 secondes évite d'exposer le jeton de session dans
  l'URL, donc dans les journaux du serveur frontal \\
\addlinespace
Validation
& adresse électronique, longueur du titre, longueur de l'identifiant
  d'opération, plafond de 64 Kio par mise à jour, corps de requête limité \\
\bottomrule
\end{tabular}
\caption{Mesures de sécurité et bonne pratique correspondante.}
\label{tab:securite}
\end{table}

\subsection{Infrastructure, déploiement et suppression}

L'infrastructure entière est décrite dans un seul modèle de pile, ce qui rend le
déploiement reproductible : une instance, une adresse élastique, un groupe de
sécurité, dix tables, un compartiment de stockage, une file avec file de
lettres mortes, un rôle d'instance et un groupe de journaux.

Le script de déploiement enchaîne six étapes, et chacune échoue franchement :

\begin{enumerate}[leftmargin=1.4em,itemsep=2pt]
  \item vérification de l'identité, de la région, et de l'\emph{éligibilité} du
        type d'instance — interrogée auprès du service, et non tenue dans une
        liste à la main ;
  \item création éventuelle de la paire de clés, puis des deux secrets ;
  \item déploiement de la pile, avec suppression préalable si elle se trouve
        dans un état d'échec antérieur, qui interdit toute mise à jour ;
  \item construction du backend puis du front, en libérant les dépendances au
        fur et à mesure pour tenir dans l'espace disque disponible ;
  \item envoi de l'archive vers le stockage, puis installation sur l'instance
        par commande à distance ;
  \item vérification que le service répond, que le parcours visiteur est
        relayé, et que le micro-cache est actif.
\end{enumerate}

La suppression mérite un script dédié, parce qu'une simple suppression de pile
échoue : le compartiment est versionné, donc il faut purger les objets
\emph{et} leurs versions ; les paramètres et la paire de clés vivent hors de la
pile et survivent à sa suppression. Le script purge, supprime, puis \emph{vérifie}
huit familles de ressources, en distinguant trois états : absent, restant, et
\emph{inconnu} — une commande de contrôle qui échoue ne doit jamais être comptée
comme « rien ne reste ».

\subsection{Limites et risques}

\begin{table}[H]
\centering
\small
\begin{tabular}{@{}p{4.4cm}p{11.2cm}@{}}
\toprule
\textbf{Limite} & \textbf{Conséquence, et raison du choix} \\
\midrule
Trafic sans chiffrement
& les jetons et les adresses circulent en clair. Un certificat exige un nom de
  domaine, absent du sujet. Limite assumée, documentée, et non dissimulée \\
\addlinespace
Instance unique
& aucune haute disponibilité : un arrêt de la machine interrompt le service.
  C'est le prix d'un écrivain unique par document, qui simplifie C2. Le
  redémarrage automatique ramène le processus, pas la machine \\
\addlinespace
Identité sans mot de passe
& quiconque connaît une adresse peut écrire en son nom. Le sujet l'autorise ;
  c'est aussi ce qui rend l'identité du visiteur inutile \\
\addlinespace
Jeton non révocable
& un jeton volé reste valable jusqu'à son expiration. Sans liste de révocation,
  il n'y a pas de contournement autre que le renouvellement \\
\addlinespace
Rattrapage plafonné
& la lecture du journal s'arrête à un nombre d'entrées fixé, sans signaler la
  troncature : un client doit boucler. Un appel unique sur un document ancien
  donnerait un journal partiel, donc une vérification de convergence fausse \\
\addlinespace
Restauration non facturée
& une restauration de version est gratuite et franchit un blocage. C'est un
  recours voulu, mais le compteur ne reflète donc pas tout le travail effectué \\
\addlinespace
Bac à sable du courriel
& tant que l'accès de production n'est pas accordé, seule l'adresse vérifiée
  reçoit la notification \\
\bottomrule
\end{tabular}
\caption{Limites connues et risques associés.}
\label{tab:limites}
\end{table}

\hypothese{L'absence de haute disponibilité est présentée comme une limite
assumée, non comme un oubli. Le cours demande explicitement de savoir et de dire
qu'une instance unique n'est pas disponible. Les pistes d'évolution — plusieurs
instances avec élection d'un écrivain par document, ou externalisation de l'état
— sont plus coûteuses que le bénéfice attendu pour un projet évalué sur la
justification des choix.}
